% §4 Exercises
%
% Cross-references and their resolution status:
%
% Already resolved (within chapter 4):
%   \Cref{ax:element-indep}      -- Axiom 1 formal statement (§4.2)
%   \Cref{ax:block-rates-indep}  -- Axiom 2 formal statement (§4.3)
%   \Cref{thm:kronecker}         -- Kronecker Factorization Theorem (§4.4)
%   \Cref{cor:parameter-reduction} -- Parameter count corollary (§4.4)
%   \Cref{def:bernoulli-bool}    -- Formal definition of Bernoulli[bool] (§4.5)
%   \Cref{ex:biased-coin-formal} -- Biased-coin worked example (§4.5)
%   \Cref{prop:stable-Q}         -- Stable channel matrices proposition (§4.3)
%
% Unresolved forward refs (will resolve in Plan 4, Part I integration):
%   \Cref{sec:and-under-noise}   -- chapter 3, section 3.1 (AND under noise)
%
% Difficulty-marker macros (not in alex.sty; defined here for chapter scope):
\newcommand{\exstar}{$\star$}
\newcommand{\exstarstar}{$\star\star$}
\newcommand{\exstarstarstar}{$\star\star\star$}

\section*{Exercises}
\addcontentsline{toc}{section}{Exercises}

\begin{enumerate}

% -----------------------------------------------------------------------
% WARM-UP  (★ × 3)
% -----------------------------------------------------------------------

\item \exstar
\textbf{Axiom 1 in measure-theoretic notation.}
Let $(\Omega, \mathcal{F}, \mathbb{P})$ be the underlying probability space,
and for each $x \in U$ define the error event $E_x = \{\omega \in \Omega :
\tilde{A}(x)(\omega) \neq \mathbf{1}_A(x)\} \in \mathcal{F}$.
Reformulate \Cref{ax:element-indep} (Axiom~1) using this notation: write out
the independence condition both pairwise and for an arbitrary finite family
$\{x_1, \ldots, x_k\} \subseteq U$.
Does the condition depend on the specific probability measure $\mathbb{P}$,
or only on the sigma-algebra $\mathcal{F}$?

\item \exstar
\textbf{Axiom 2 as a constraint on $Q$.}
Restate \Cref{ax:block-rates-indep} (Axiom~2) as a property of the channel
matrix $Q$. Specifically: in the order-2 Boolean case, $Q$ has two rows.
Axiom~2 declares the random quantities governing those rows to be
conditionally independent. Express this as a statement about the joint
density $f(\fprate, \fnrate \Given A)$.
What behavior does Axiom~2 \emph{forbid} $Q$ from exhibiting across a
sequence of queries?
(See \Cref{prop:stable-Q} for the stability consequence derived in §4.3.)

\item \exstar
\textbf{Computing a Kronecker product.}
Let
\[
  A \;=\; \begin{pmatrix} 1 & 2 \\ 3 & 4 \end{pmatrix},
  \qquad
  B \;=\; \begin{pmatrix} 5 & 6 \\ 7 & 8 \end{pmatrix}.
\]
Compute $A \otimes B$ by hand. Verify that $A \otimes B$ is a $4 \times 4$
matrix. Then compute $B \otimes A$ and check whether
$A \otimes B = B \otimes A$.
(\textit{Hint:} the Kronecker product is not in general commutative;
the two products will differ, but both will be $4 \times 4$.)

% -----------------------------------------------------------------------
% MEDIUM  (★★ × 5)
% -----------------------------------------------------------------------

\item \exstarstar
\textbf{Proving the parameter-count corollary.}
Prove \Cref{cor:parameter-reduction} by direct counting.
\begin{enumerate}
  \item[(a)] Without any independence assumption, an $n$-th order model on
    $m$ joint queries has $n(2^m - 1)$ free parameters. Explain where this
    count comes from.
  \item[(b)] Under Axioms~1 and~2, the Kronecker factorization
    (\Cref{thm:kronecker}) forces the parameter count down to $nm$.
    Show why: how many free parameters does each per-query matrix
    $Q_i$ contribute, and why do the $m$ per-query matrices together
    total $nm$?
  \item[(c)] For the order-2 Boolean channel ($n = 2$), compute the ratio
    $n(2^m - 1) / (nm)$ for $m = 10$ and $m = 50$. Comment on the
    magnitude of the reduction.
\end{enumerate}

\item \exstarstar
\textbf{Verifying the Kronecker theorem by hand for $m = 2$.}
Use the asymmetric channel from the biased-coin worked example
(\Cref{ex:biased-coin-formal}) with $\fprate = 0.05$ and $\fnrate = 0.10$,
so $Q = \bigl(\begin{smallmatrix} 0.90 & 0.10 \\ 0.05 & 0.95 \end{smallmatrix}\bigr)$.
\begin{enumerate}
  \item[(a)] Compute $Q \otimes Q$ explicitly. Write out all sixteen entries
    of the resulting $4 \times 4$ matrix.
  \item[(b)] For the two-query input pair $(\top, \bot)$, use direct
    probability to compute $\Prob{\tilde{A}(\top) = s_1,\, \tilde{A}(\bot) = s_2
    \Given A}$ for all four outcome pairs $(s_1, s_2) \in \{\top,\bot\}^2$.
  \item[(c)] Verify that the entries from (b) match the corresponding row of
    $Q \otimes Q$ from (a), confirming \Cref{thm:kronecker} for $m = 2$.
\end{enumerate}

\item \exstarstar
\textbf{A Markov channel violates Axiom~1.}
The Markov-correlated channel in §4.3 couples the error rate on each query
to the outcome of the previous query: if the previous query produced an error,
the current query uses elevated false-positive rate $\fprate_1 > \fprate_0$.
Consider two consecutive queries on distinct elements $x, y \in B_2 = \{\bot\}$,
with base rate $\fprate_0 = 0.05$ and elevated rate $\fprate_1 = 0.30$.
\begin{enumerate}
  \item[(a)] Compute $\Prob{E_x = 1,\, E_y = 1}$ under the Markov channel,
    where the channel starts in the correct state.
  \item[(b)] Compute $\Prob{E_x = 1} \cdot \Prob{E_y = 1}$ under the same
    parameters.
  \item[(c)] Show that these two quantities differ, thereby exhibiting a
    violation of \Cref{ax:element-indep} (Axiom~1).
\end{enumerate}
(\textit{Note:} §4.3 frames this example as an Axiom~2 failure; here we recast
it explicitly as an Axiom~1 failure. Both framings are valid; the Markov
coupling violates both axioms simultaneously.)

\item \exstarstar
\textbf{The implicit Axiom-2 invocation in §3.1's AND derivation.}
Re-read the derivation of AND under noise in \Cref{sec:and-under-noise}.
The derivation writes the joint probability that both inputs to an AND gate
produce errors as a product of two individual error probabilities.
\begin{enumerate}
  \item[(a)] Identify the step where \Cref{ax:block-rates-indep} (Axiom~2)
    is invoked implicitly. Quote or describe the exact expression where
    the factoring occurs.
  \item[(b)] Suppose Axiom~2 failed: the false-positive rate on the first
    input told you something about the false-positive rate on the second.
    Write an expression for the joint error probability in that case
    (in terms of the conditional density $f(\fprate_2 \Given \fprate_1)$).
  \item[(c)] Would the closed-form AND output rate from §3.1 survive if
    Axiom~2 failed? Explain in one or two sentences.
\end{enumerate}

\item \exstarstar
\textbf{Constructing a $\Bernoulli{\BoolType}$ from a biased coin.}
Following \Cref{def:bernoulli-bool} and \Cref{ex:biased-coin-formal},
construct a $\Bernoulli{\BoolType}$ instance for a biased coin with
bias $p = 0.7$, false-positive rate $\fprate = 0.10$, and false-negative
rate $\fnrate = 0.05$.
\begin{enumerate}
  \item[(a)] Write down the channel matrix $Q$ and the joint PMF over all
    four (true value, observed report) pairs.
  \item[(b)] Verify Axiom~1 (\Cref{ax:element-indep}) in the following sense:
    the Boolean universe has only two elements ($\top$ and $\bot$), so state
    what element-wise independence means in this degenerate case and confirm
    it holds by construction.
  \item[(c)] Verify Axiom~2 (\Cref{ax:block-rates-indep}) in the following
    sense: state the joint density condition $f(\fprate, \fnrate \Given A) =
    f(\fprate \Given A)\,f(\fnrate \Given A)$ and explain why it is satisfied
    when the observation mechanism acts independently on each partition block.
\end{enumerate}

% -----------------------------------------------------------------------
% CHALLENGING / EXTENSION  (★★★ × 3)
% -----------------------------------------------------------------------

\item \exstarstarstar
\textbf{Counterexample: Axiom~1 does not imply Axiom~2.}
Construct a probability model that satisfies \Cref{ax:element-indep}
(Axiom~1: error events at distinct elements are mutually independent) but
violates \Cref{ax:block-rates-indep} (Axiom~2: per-block error rates are
conditionally independent across blocks).
\begin{enumerate}
  \item[(a)] Describe the model explicitly: specify $U$, the partition blocks,
    and the joint distribution of error rates across blocks.
  \item[(b)] Verify that Axiom~1 holds: for any fixed realization of the
    block rates, the per-element error events are independent.
  \item[(c)] Show that Axiom~2 fails: the joint density of the block rates
    does not factor.
\end{enumerate}
(\textit{Hint:} consider a two-block model where the rates $\fprate$ and
$\fnrate$ are drawn from a bivariate distribution with nonzero correlation,
such as a Dirichlet or a copula, while the per-element draws given those
rates remain independent.)

\item \exstarstarstar
\textbf{Without Axiom~2, the channel-matrix view collapses.}
Show that without \Cref{ax:block-rates-indep} (Axiom~2), the joint
distribution over $m$ queries to an element in a fixed block $B_i$ cannot
in general be written as a product of $m$ identical channel-matrix rows.
\begin{enumerate}
  \item[(a)] Give a concrete two-query example where the marginal per-query
    distribution is consistent with a channel matrix $Q$, yet the
    \emph{joint} two-query distribution over outcomes for a single block
    $B_i$ does not equal $(Q_{i,\cdot})^{\otimes 2}$.
    (You may use the Markov-correlated channel from §4.3 as your example.)
  \item[(b)] What mathematical object would replace the single matrix $Q$
    in a model without Axiom~2? Name or describe an appropriate
    generalization (e.g., a time-varying sequence of matrices, a hidden
    Markov model, or a conditional channel).
  \item[(c)] Explain in one paragraph why the Kronecker factorization theorem
    (\Cref{thm:kronecker}) depends on Axiom~2 and not only on Axiom~1.
\end{enumerate}

\item \exstarstarstar
\textbf{Tightness of the parameter-count bound.}
\Cref{cor:parameter-reduction} states that under Axioms~1 and~2, an
$n$-th order Bernoulli predicate on $m$ queries requires $nm$ free
parameters. Argue that this bound is tight: there exist $\Bernoulli{\BoolType}$
instances satisfying both axioms that \emph{genuinely} require $2m$ parameters
to specify (in the order-2 Boolean case).
\begin{enumerate}
  \item[(a)] Construct a family of $\Bernoulli{\BoolType}$ instances
    parameterized by $m$ independently chosen channel matrices
    $Q_1, Q_2, \ldots, Q_m$, each with its own $(\fprate_i, \fnrate_i)$.
    Verify that this family satisfies both axioms.
  \item[(b)] Show that no two distinct members of the family yield the same
    joint distribution over $m$ queries, so distinct parameter tuples
    $(\fprate_1, \fnrate_1, \ldots, \fprate_m, \fnrate_m)$ produce distinct
    joint confusion matrices.
  \item[(c)] Conclude that the $2m$ parameters are not merely sufficient but
    necessary: no reparameterization can reduce the family to fewer than $2m$
    free parameters without collapsing distinct members.
\end{enumerate}

\end{enumerate}
